\documentclass[11pt]{article}
\usepackage[a4paper,margin=1in]{geometry}
\usepackage{fontspec}
\usepackage[boldfont=false]{xeCJK}
\setCJKmainfont{FandolSong-Regular.otf}
\usepackage{amsmath}
\usepackage{amssymb}
\usepackage{array}
\usepackage{booktabs}
\usepackage{graphicx}
\usepackage{adjustbox}
\usepackage{enumitem}
\setlist[itemize]{topsep=2pt,partopsep=0pt,parsep=0pt,itemsep=2pt,leftmargin=1.4em}
\setlist[enumerate]{topsep=2pt,partopsep=0pt,parsep=0pt,itemsep=2pt,leftmargin=1.6em}
\usepackage{parskip}
\usepackage{fvextra}
\fvset{breaklines=true,breakanywhere=true,fontsize=\small}
\setlength{\parindent}{0pt}

\begin{document}

\begin{center}
  {\LARGE\bfseries Human resource management (A Level)}\\[0.35em]
  {\large A Level Business}
\end{center}
\vspace{0.6em}

\section*{Organisational structure}

An \textbf{organisational structure} 组织结构 shows how a business arranges its people: who reports to whom, and who does what. A good structure makes roles clear and helps the firm meet its objectives. It is often drawn as an organisation chart.

\section*{Key terms in structure}

\begin{itemize}
\item a \textbf{hierarchy} 层级 is the set of levels in a firm, from the top managers down to the shop-floor workers. A "tall" structure has many levels; a "flat" structure has few.

\item the \textbf{chain of command} 指挥链 is the line along which orders pass down, from the top to the bottom.

\item the \textbf{span of control} 管理幅度 is the number of staff one manager is directly in charge of. A wide span means many staff per manager; a narrow span means few.

\end{itemize}

A tall structure has narrow spans and a long chain of command, so messages travel slowly. A flat structure has wide spans, so managers must trust staff more.

\section*{Centralisation and decentralisation}

\begin{itemize}
\item \textbf{centralisation} 集权 — most decisions are made at the top. This gives strong control and quick, consistent decisions, but ignores local knowledge.

\item \textbf{decentralisation} 分权 — decisions are shared out to lower levels and local branches. This speeds up decisions and motivates staff, but control is weaker.

\end{itemize}

Most firms sit somewhere between the two.

\section*{Delayering and matrix structures}

\textbf{Delayering} 精简层级 means removing one or more levels of management. It cuts cost and shortens the chain of command, but it gives the remaining managers more to do.

A \textbf{matrix structure} 矩阵结构 groups staff by both their department and the project they work on, so a worker may report to two managers. It is good for teamwork across departments, but the two-boss system can cause confusion.

\section*{Delegation and accountability}

\textbf{Delegation} 授权 means a manager gives a task, and the power to do it, to a more junior worker. It frees the manager's time, develops staff, and motivates them.

But a manager who delegates is still \textbf{accountable} — the \textbf{accountability} 问责 for the result stays with them. So managers must delegate to people they trust, and still check the outcome.

\section*{Business communication}

\textbf{Communication} 沟通 is the passing of a message from one person to another. Effective communication follows a clear path:

\begin{itemize}
\item a \textbf{sender} 发送者 has a message.

\item the message is sent through a \textbf{medium} 媒介 (e.g. a meeting, email or notice).

\item a \textbf{receiver} 接收者 takes in the message.

\item \textbf{feedback} 反馈 shows the message was understood.

\end{itemize}

Good communication raises motivation, cuts mistakes, and speeds up decisions. Methods can be verbal (spoken), written, or visual; the best method depends on the message, the cost, and the need for a record.

\section*{Barriers to communication}

A \textbf{barrier} 障碍 is anything that stops a message getting through clearly. Common barriers are:

\begin{itemize}
\item a message that is too long or unclear.

\item a poor medium (e.g. a noisy phone line).

\item too many levels in the hierarchy, so the message is changed on the way.

\item language or cultural differences, and lack of trust.

\end{itemize}

To improve communication, a firm can shorten the chain of command, choose the right medium, keep messages short and clear, and always ask for feedback.

\section*{Leadership styles}

A \textbf{leadership style} 领导风格 is the way a \textbf{leader} 领导 makes decisions and treats staff. Four styles are:

\begin{center}
\adjustbox{max width=\linewidth}{%
\begin{tabular}{ll}
\toprule
\textbf{Style} & \textbf{How decisions are made} \\
\midrule
\textbf{autocratic} 专制型 & the leader decides alone and tells staff what to do \\
\textbf{democratic} 民主型 & the leader asks staff for ideas before deciding \\
\textbf{paternalistic} 家长式 & the leader decides, but in what they see as the staff's best interests \\
\textbf{laissez-faire} 放任型 & the leader gives staff freedom to decide for themselves \\
\bottomrule
\end{tabular}}
\end{center}

No style is always best. The right one depends on the task, the staff, and the time available — an urgent crisis may need an autocratic style, while skilled staff may work best under a democratic or laissez-faire one.

\section*{Emotional intelligence and leading change}

\textbf{Emotional intelligence} 情商 is a leader's skill in understanding their own and other people's feelings, and managing them well. Leaders with high emotional intelligence build trust and handle conflict better.

This matters most when \textbf{leading change}. People often resist change because they fear losing their job or status. A good leader explains why the change is needed, involves staff, and supports them through it.

\section*{Hard and soft HRM}

There are two broad approaches to \textbf{human resource management} 人力资源管理 (HRM):

\begin{itemize}
\item \textbf{hard HRM} 硬性人力资源管理 treats staff as a resource like any other — to be used at the lowest cost. It uses tight control and short-term contracts.

\item \textbf{soft HRM} 软性人力资源管理 treats staff as the firm's most valuable asset — to be developed and motivated. It uses training, involvement and long-term careers.

\end{itemize}

Both link to \textbf{workforce planning} 人力规划: working out the staff the firm will need, and how to get them.

\section*{Employer-employee relations and trade unions}

The \textbf{management of change} 变革管理 works best when \textbf{employer-employee relations} 劳资关系 (the link between bosses and workers) are good.

A \textbf{trade union} 工会 is an organised group of workers that speaks for its members. Through \textbf{collective bargaining} 集体谈判, the union and the employer negotiate pay and conditions for the whole group, instead of one worker at a time. If talks fail, workers may take \textbf{industrial action} 劳工行动, such as a \textbf{strike} 罢工 (stopping work). Good relations and fair negotiation avoid this and keep the business running.


\end{document}
